\documentclass[12pt]{article}
\usepackage[a4paper,margin=1in]{geometry}
\usepackage[utf8]{inputenc}
\usepackage[T1]{fontenc}
\usepackage{amsmath,amssymb}
\usepackage{textcomp}
\usepackage{parskip}
\begin{document}
\begin{center}
{\Large \textbf{2017 STEP I}}
\end{center}

\subsection*{Page 1}
\begin{flushleft}
\bigskip\textbf{Section A: Pure Mathematics}\par
\bigskip\textbf{Question 1.} (i) Use the substitution u = x sin x + cos x to find\par
x\par
dx .\par
\[\leq x tan x + 1\]
Find by means of a similar substitution, or otherwise,\par
x\par
dx .\par
\[\leq x cot x 1\]
-\par
\textbf{(ii)} Use a substitution to find\par
\[x sec2 x tan x\]
dx\par
\[\leq x sec2 x tan x\]
-\par
and\par
x sin x cos x\par
dx .\par
\[\leq (x sin x cos x)2\]
-\par
\bigskip\textbf{Question 2.} (i) The inequality \leq 1 holds for t \geq 1. By integrating both sides of this inequality over\par
t\par
\[the interval 1 \leq t \leq x, show that\]
\[ln x \leq x 1 ( )\]
- ∗\par
\[for x \geq 1. Show similarly that ( ) also holds for 0 < x \leq 1.\]
∗\par
1 1\par
\textbf{(ii)} Starting from the inequality \leq for t \geq 1, show that\par
\[t^2 t\]
\[ln x \geq 1 ( )\]
- x ∗∗\par
\[for x > 0.\]
\textbf{(iii)} Show, by integrating ( ) and ( ), that\par
∗ ∗∗\par
2 ln y y + 1\par
\[\leq \leq\]
y + 1 y 1 2y\par
-\par
\[for y > 0 and y = 1.\]
\[\geq\]
\[\geq\geq\theta\mu\nu) \phi \geq\]
\end{flushleft}
\subsection*{Page 2}
\begin{flushleft}
\bigskip\textbf{Question 3.} The points P (ap2, 2ap) and Q(aq2, 2aq), where p > 0 and q < 0, lie on the curve C with\par
equation\par
\[y = 4ax ,\]
\[where a > 0 . Show that the equation of the tangent to C at P is\]
\[y = x + ap .\]
p\par
The tangents to the curve at P and at Q meet at R. These tangents meet the y-axis at S\par
and T respectively, and O is the origin. Prove that the area of triangle OP Q is twice the\par
area of triangle RST .\par
\bigskip\textbf{Question 4.} (i) Let r be a real number with r < 1 and let\par
| |\par
\infty\par
n\par
\[S = r .\]
\[n\ne0\]
\[You may assume without proof that S = .\]
1 r\par
-\par
\[Let p = 1 + r + r^2. Sketch the graph of the function 1 + r + r^2 and deduce that\]
\[3 \leq p < 3 .\]
\[Show that, if 1 < p < 3, then the value of p determines r, and hence S, uniquely.\]
\[Show also that, if 3 < p < 1, then there are two possible values of S and these values\]
\[satisfy the equation (3 p)S^2 3S + 1 = 0.\]
- -\par
\textbf{(ii)} Let r be a real number with r < 1 and let\par
| |\par
\infty\par
n-1\par
\[T = nr .\]
\[n\ne1\]
\[You may assume without proof that T = .\]
(1 r)2\par
-\par
\[Let q = 1 + 2r + 3r2. Find the set of values of q that determine T uniquely.\]
Find the set of values of q for which T has two possible values. Find also a quadratic\par
equation, with coefficients depending on q, that is satisfied by these two values.\par
\[\geq\geq\theta\mu\nu) \pi\]
\end{flushleft}
\subsection*{Page 3}
\begin{flushleft}
\bigskip\textbf{Question 5.} A circle of radius a is centred at the origin O. A rectangle P QRS lies in the minor sector\par
OM N of this circle where M is (a, 0) and N is (a cos β, a sin β), and β is a constant with\par
\[\pi\]
0 < β < . Vertex P lies on the positive x-axis at (x, 0); vertex Q lies on ON ; vertex R lies\par
on the arc of the circle between M and N ; and vertex S lies on the positive x-axis at (s, 0).\par
Show that the area A of the rectangle can be written in the form\par
\[A = x(s x) tan β .\]
-\par
Obtain an expression for s in terms of a, x and β, and use it to show that\par
\[dA x^2\]
\[= (s 2x) tan β tan β .\]
dx - - s\par
\[Deduce that the greatest possible area of rectangle P QRS occurs when s = x(1 + sec β) and\]
\[show that this greatest area is 1 a^2 tan 1 β .\]
2 2\par
\[Show also that this greatest area occurs when ∠ROS = 1 β .\]
\bigskip\textbf{Question 6.} In this question, you may assume that, if a continuous function takes both positive and\par
negative values in an interval, then it takes the value 0 at some point in that interval.\par
\textbf{(i)} The function f is continuous and f(x) is non-zero for some value of x in the interval\par
\[0 \leq x \leq 1. Prove by contradiction, or otherwise, that if\]
\[f(x)dx = 0 ,\]
\[\leq\]
\[then f(x) takes both positive and negative values in the interval 0 \leq x \leq 1.\]
\textbf{(ii)} The function g is continuous and\par
1 1 1\par
2 2\par
\[g(x) dx = 1 , xg(x) dx = α , x g(x) dx = α . ( )\]
\[\leq \leq \leq ∗\]
0 0 0\par
Show, by considering\par
(x α) g(x) dx ,\par
\[\leq -\]
\[that g(x) = 0 for some value of x in the interval 0 \leq x \leq 1.\]
\[Find a function of the form g(x) = a + bx that satisfies the conditions ( ) and verify\]
∗\par
\[that g(x) = 0 for some value of x in the interval 0 \leq x \leq 1.\]
\[\geq\]
\textbf{(iii)} The function h has a continuous derivative h and\par
1 1\par
\[h(0) = 0 , h(1) = 1 , h(x) dx = β , xh(x) dx = 1 β(2 β) .\]
\[\leq \leq 2 -\]
0 0\par
\[\geq\]
\[Use the result in part (ii) to show that h (x) = 0 for some value of x in the interval\]
\[0 \leq x \leq 1.\]
\[\geq\geq\theta\mu\nu) \sigma \geq\]
\end{flushleft}
\subsection*{Page 4}
\begin{flushleft}
\bigskip\textbf{Question 7.} The triangle ABC has side lengths BC = a, CA = b and AB = c. Equilateral triangles\par
| | | | | |\par
BXC, CYA and AZB are erected on the sides of the triangle ABC, with X on the other\par
side of BC from A, and similarly for Y and Z. Points L, M and N are the centres of\par
rotational symmetry of triangles BXC, CYA and AZB respectively.\par
b\par
\textbf{(i)} Show that CM = and write down the corresponding expression for CL .\par
\[| | \sqrt\{3\} | |\]
\textbf{(ii)} Use the cosine rule to show that\par
\[6 LM 2 = a 2 + b 2 + c 2 + 4\sqrt\{3\} Δ ,\]
| |\par
where Δ is the area of triangle ABC. Deduce that LM N is an equilateral triangle.\par
Show further that the areas of triangles LM N and ABC are equal if and only if\par
\[a 2 + b 2 + c 2 = 4\sqrt\{3\} Δ .\]
\textbf{(iii)} Show that the conditions\par
2 \textasciicircum{}\circ\par
\[(a b) = 2ab 1 cos(C 60 )\]
- - - -\par
( )\par
and\par
\[a 2 + b 2 + c 2 = 4\sqrt\{3\} Δ\]
are equivalent.\par
Deduce that the areas of triangles LM N and ABC are equal if and only if ABC is\par
equilateral.\par
\bigskip\textbf{Question 8.} Two sequences are defined by a 1 = 1 and b 1 = 2 and, for n \geq 1,\par
\[an+1 = an + 2bn ,\]
\[bn+1 = 2an + 5bn .\]
\[Prove by induction that, for all n \geq 1,\]
2 2\par
\[a n + 2anbn b n = 1 . ( )\]
- ∗\par
\textbf{(i)} Let cn =\par
an\par
\[. Show that bn \geq 2 5 n-1 and use ( ) to show that\]
\[bn \times ∗\]
\[cn \sqrt\{2\} 1 as n .\]
\to - \to \infty\par
\textbf{(ii)} Show also that cn > \sqrt2 1 and hence that < \sqrt2 < cn + 1 .\par
- cn + 1\par
140 99\par
\[Deduce that < \sqrt\{2\} < .\]
99 70\par
\[\geq\geq\theta\mu\nu) \omega\]
\end{flushleft}
\subsection*{Page 5}
\begin{flushleft}
\bigskip\textbf{Section B: Mechanics}\par
\bigskip\textbf{Question 9.} A particle is projected at speed u from a point O on a horizontal plane. It passes through a\par
fixed point P which is at a horizontal distance d from O and at a height d tan β above the\par
plane, where d > 0 and β is an acute angle. The angle of projection α is chosen so that u is\par
as small as possible.\par
\textbf{(i)} Show that u2 = gd tan α and 2α = β + 90 \textasciicircum{}\circ .\par
\textbf{(ii)} At what angle to the horizontal is the particle travelling when it passes through P ?\par
Express your answer in terms of α in its simplest form.\par
10 Particles P 1, P 2, . . . are at rest on the x-axis, and the x-coordinate of Pn is n. The mass\par
n\par
of Pn is λ m. Particle P , of mass m, is projected from the origin at speed u towards P 1.\par
A series of collisions takes place, and the coefficient of restitution at each collision is e, where\par
\[0 < e < 1. The speed of Pn immediately after its first collision is un and the speed of Pn\]
immediately after its second collision is vn. No external forces act on the particles.\par
1 + e\par
\textbf{(i)} Show that u 1 = u and find expressions for un and vn in terms of e, λ, u and n.\par
1 + λ\par
\textbf{(ii)} Show that, if e > λ, then each particle (except P ) is involved in exactly two collisions.\par
\textbf{(iii)} Describe what happens if e = λ and show that, in this case, the fraction of the initial\par
kinetic energy lost approaches e as the number of collisions increases.\par
\textbf{(iv)} Describe what happens if λe = 1. What fraction of the initial kinetic energy is\par
eventually lost in this case?\par
\[\geq\geq\theta\mu\nu) ) \geq\]
\end{flushleft}
\subsection*{Page 6}
\begin{flushleft}
\bigskip\textbf{Question 11.} A plane makes an acute angle α with the horizontal. A box in the shape of a cube is fixed\par
onto the plane in such a way that four of its edges are horizontal and two of its sides are\par
vertical.\par
A uniform rod of length 2L and weight W rests with its lower end at A on the bottom of the\par
box and its upper end at B on a side of the box, as shown in the diagram below. The vertical\par
plane containing the rod is parallel to the vertical sides of the box and cuts the lowest edge\par
of the box at O. The rod makes an acute angle β with the side of the box at B.\par
The coefficients of friction between the rod and the box at the two points of contact are both\par
\[tan γ, where 0 < γ < 1 \pi.\]
The rod is in limiting equilibrium, with the end at A on the point of slipping in the direction\par
\[away from O and the end at B on the point of slipping towards O. Given that α < β, show\]
\[that β = α + 2γ.\]
[Hint: You may find it helpful to take moments about the midpoint of the rod.]\par
B\par
β\par
A\par
α\par
O\par
\[\geq\geq\theta\mu\nu) (\]
\end{flushleft}
\subsection*{Page 7}
\begin{flushleft}
\bigskip\textbf{Section C: Probability and Statistics}\par
\bigskip\textbf{Question 12.} In a lottery, each of the N participants pays \pounds c to the organiser and picks a number from\par
1 to N . The organiser picks at random the winning number from 1 to N and all those\par
participants who picked this number receive an equal share of the prize, \pounds J .\par
\textbf{(i)} The participants pick their numbers independently and with equal probability. Obtain\par
an expression for the probability that no participant picks the winning number, and\par
hence determine the organiser's expected profit.\par
Use the approximation\par
a N\par
-a\par
1 e ( )\par
\[\int - N \sum \approx ∗\]
\[to show that if 2N c = J then the organiser will expect to make a loss.\]
\[Note: e > 2.\]
\textbf{(ii)} Instead of the numbers being equally popular, a fraction γ of the numbers are popular\par
and the rest are unpopular. For each participant, the probability of picking any given\par
a b\par
popular number is and the probability of picking any given unpopular number is .\par
N N\par
Find a relationship between a, b and γ.\par
Show that, using the approximation ( ), the organiser's expected profit can be expressed\par
∗\par
in the form\par
-a -b\par
Ae + Be + C ,\par
where A, B and C can be written in terms of J , c, N and γ.\par
\[In the case γ = 1 and a = 9b, find a and b. Show that, if 2N c = J , then the organiser\]
will expect to make a profit.\par
\[Note: e < 3.\]
\[\geq\geq\theta\mu\nu) ) \geq\]
\end{flushleft}
\subsection*{Page 8}
\begin{flushleft}
\bigskip\textbf{Question 13.} I have a sliced loaf which initially contains n slices of bread. Each time I finish setting a\par
STEP question, I make myself a snack: either toast, using one slice of bread; or a sandwich,\par
using two slices of bread. I make toast with probability p and I make a sandwich with\par
probability q, where p + q = 1, unless there is only one slice left in which case I must, of\par
course, make toast.\par
Let sr (1 \leq r \leq n) be the probability that the rth slice of bread is the second of two slices\par
used to make a sandwich and let tr (1 \leq r \leq n) be the probability that the rth slice of bread\par
is used to make toast. What is the value of s 1?\par
Explain why the following equations hold:\par
\[tr = (sr-1 + tr-1) p (2 \leq r \leq n 1) ;\]
-\par
\[sr = 1 (sr-1 + tr-1) (2 \leq r \leq n) .\]
-\par
\[Hence, or otherwise, show that sr = q(1 sr-1) for 2 \leq r \leq n 1 .\]
- -\par
Show further that\par
r\par
q + ( q)\par
\[sr = - (1 \leq r \leq n 1) ,\]
1 + q -\par
and find the corresponding expression for tr.\par
Find also expressions for sn and tn in terms of q.\par
\[\geq\geq\theta\mu\nu) \nu\mu\]
\end{flushleft}
\end{document}
